% Part: history
% Chapter: biographies
% Section: bertrand-russell

\documentclass[../../../include/open-logic-section]{subfiles}

\begin{document}

\olfileid[pt]{his}{bio}{rus} 

\olsection{Bertrand Russell}

\olphoto{russell-bertrand}{Bertrand Russell}

Bertrand Russell é celebrado como um dos fundadores da filosofia
analítica moderna. Nascido em 18 de maio de 1872, Russell não era
conhecido apenas pelo trabalho em filosofia e lógica, mas escreveu
muitos livros de divulgação em diversas áreas. Foi também um ardente
ativista político ao longo da vida.

Russell nasceu em Trellech, Monmouthshire, no País de Gales. Os pais
pertenciam à nobreza britânica. Eram livre-pensadores e chegaram a
travar amizade com os radicais de Boston da época.  Infelizmente, os
pais de Russell morreram quando ele era ainda criança, e Russell foi
viver com os avós. Ali recebeu uma educação religiosa (algo que os
pais tinham querido evitar a todo custo). A avó era muito rigorosa
em todas as questões de moralidade. Na adolescência foi sobretudo
educado em casa por professores particulares.

A influência de Russell na filosofia analítica, e sobretudo na lógica,
é enorme. Estudou matemática e filosofia no Trinity College, em
Cambridge, onde foi influenciado pelo matemático e filósofo Alfred
North Whitehead.  Em 1910, Russell e Whitehead publicaram o primeiro
volume de \emph{Principia Mathematica}, defendendo a ideia de que a
matemática é redutível à lógica. Seguiu-se a publicação de centenas de
livros, ensaios e panfletos políticos. Em 1950, ganhou o Prêmio Nobel
de Literatura.

Russell se encontrava profundamente envolvido na política e no
ativismo social. Durante a Primeira Guerra Mundial foi preso e
condenado a seis meses de prisão por atividades pacifistas e
protesto. Na prisão pôde escrever e ler, e afirma ter achado a
experiência ``bastante agradável.'' Permaneceu pacifista ao longo da
vida, e voltou a ser detido por participar num comício pelo
desarmamento nuclear em 1961. Sobreviveu também a um acidente de avião
em 1948, no qual os únicos sobreviventes eram os que estavam sentados
na área de fumantes. Por isso, Russell declarou que devia a vida ao
tabagismo. Russell se casou quatro vezes, mas tinha fama de manter
relações extraconjugais.  Morreu em 2 de fevereiro de 1970, aos 97
anos, em Penrhyndeudraeth, no País de Gales.

\begin{reading}
Russell escreveu uma autobiografia em três partes, que abrange sua
vida de 1872 a 1967 \citep{Russell1967,Russell1968,Russell1969}.  O
Bertrand Russell Research Centre na McMaster University alberga os
arquivos de Bertrand Russell. Ver o site em \citet{Duncan2015},
para informação sobre os volumes de suas obras reunidas (incluindo
índices pesquisáveis) e projetos de arquivo.  O artigo de Russell
\emph{On Denoting} \citep{Russell1905} é um clássico da filosofia
analítica do século XX.

A entrada da Stanford Encyclopedia of Philosophy sobre Russell
\citep{Irvine2015} inclui excertos de áudio de Russell a falar sobre
Desejo e teoria política. Estão disponíveis online muitas entrevistas
em vídeo com Russell. Para o ouvir falar de tabagismo e do envolvimento
num acidente de avião, ver, por exemplo, \citet{RussellND}. Algumas
obras de Russell, incluindo a \emph{Introduction to Mathematical
Philosophy}, estão disponíveis como audiolivros gratuitos em
\citet{LibriVoxND}.
\end{reading}

\end{document}
